\section{Related Work}
This chapter provides a quick overview of the literature on unsupervised learning approaches for point clouds as well as variety of local feature extraction techniques on point clouds.

\vspace{5mm}

Since deep learning methods can directly consume point clouds, feature learning for 3D objects has advanced significantly in the recent years. By independently learning every point in the point cloud and combining point characteristics with permutation invariant operations, PointNet \cite{qi2017pointnet} and Deep Sets \cite{zaheer2017deep} were the pioneers in the field to handle unordered and unstructured 3D points. Both PointNet\cite{qi2017pointnet} and Deep Sets\cite{zaheer2017deep} were important developments in deep learning for handling unordered data; PointNet was designed particularly for 3D point cloud analysis, while Deep Sets provides a broader framework for feature learning that is permutation invariant and could be extrapolated to graphs and sequences. Although these methods have shown promising results in a number of applications, further research is required to resolve their shortcomings in capturing local features and enhance their scalability and performance in a variety of applications and larger datasets. By adding a hierarchical feature learning mechanism that operates directly in a metric space, PointNet++ \cite{qi2017pointnet++} expands the capabilities of PointNet. Its contributions include the ability to learn features at different scales, the hierarchical organization of point clouds into nested partitions and sampling/grouping procedures that record intricate spatial relationships while preserving permutation invariance. But because of the increased computational complexity brought forth by hierarchical architecture, training and inference periods may be prolonged. Its applicability in some circumstances may be limited by its scalability, especially with very large point clouds. PointNet++ improves point cloud feature learning overall, but its practical application necessitates careful consideration of processing resources and scalability.

\vspace{5mm}

\cite{wang2019dynamic} designed a \ac{DGCNN} architecture specifically for point cloud data. Its main contribution is to enable more efficient feature learning by dynamically building graph structures from point clouds based on local geometric correlations. But on the other hand, dynamic graph creation adds more computational complexity and might necessitate precise hyperparameter tweaking. While PointCNN\cite{li2018pointcnn} achieved state-of-the-art performance on various point cloud tasks, the X-transformation introduces additional computational overhead and may limit scalability, specially with large-scale datasets. PointCNN\cite{li2018pointcnn} proposes a \ac{CNN} specifically designed for processing point clouds. Its major contribution is the introduction of X-transformation, a permutation-invariant operation that reorients point clouds into a canonical layout, allowing application of standard convolutional operations. PointConv \cite{wu2019pointconv}, a deep learning architecture that conducts convolutional operations directly on 3D point clouds was presented in "Deep convolutional Networks on 3D Point Clouds" \cite{wu2019deep}. Its contribution is the construction of a new convolutional operator that is permutation invariant and preserves local geometric structures of point clouds. However, the convolutional procedures on irregular point clouds may cause PointConv to suffer from increased computational complexity, which could affect the efficiency of training and inference.

\vspace{5mm}

\cite{liu2019relation} presented a \ac{CNN} that was specifically designed to handle tasks related to point cloud analysis. The incorporation of relation-aware convolutional procedures, which capture both local and global geometric relationships inside point clouds, is its primary contribution. To obtain optimal performance, careful parameter tuning may be necessary due to the dependence on relation-aware features, which may result in an increase in computing overhead. Position-adaptive feature learning was the primary innovation of PAConv\cite{xu2021paconv}, which is the dynamic assembly of convolutional kernels based on local geometric structures. By using this method, the network performed better in segmentation and classification tasks by improving its capacity to detect intricate geometric patterns and variances within point clouds. Nevertheless, computational complexity may provide PAConv with difficulties, especially during training and inference. More research is needed to determine whether or not it can scale to large-size datasets. All things considered, PAConv was a major step forward in point cloud feature learning, but its practical application necessitates careful attention to scalability and processing resources. The main contribution of AdaptConv \cite{zhou2021adaptive} was the inclusion of adaptive graph convolutional operations, which dynamically modify convolutional kernels receptive fields according to the point clouds local geometric characteristics. Because of the flexibility, the network can better execute tasks like segmentation and classification by capturing local and global geometric data. Convolutional operations adaptive nature might add more computing complexity and necessitate precise parameter tweaking. Moreover, the quality of local geometric information and intricacy of the point cloud may have a significant influence on the adaptive graph convolutional approach's efficacy, which may affect how well it performs in various datasets and applications. 

\vspace{5mm}

A unique method for encoding point clouds in 3D space was proposed in \cite{sun2021quadratic}, which uses quadratic terms to more precisely reflect local surface geometry. The main contribution is the addition of quadratic components to the conventional point-to-point representation, which improves the model's ability to detect surface curvature and form fluctuations. The suggested approach enhances the discriminative power of deep learning models for tasks like point cloud segmentation and classification by adding higher-order geometric information to the representation. Quadratic terms, on the other hand, might add more computational complexity, which could affect the speed of inference and training. Furthermore, the quality and density of raw point cloud data may have an impact on how successful the suggested representation is, necessitating careful thought and preprocessing. All things considered, \cite{sun2021quadratic} presented improvements in encoding intricate geometric details. However, due to its computational complexity and data sensitivity issues, more research is needed before it can be used in real-world scenarios. \cite{ding2021point} described how perturbation learning can be used to improve point cloud data resolution. This method's primary advantage is its capacity to produce high resolution point clouds with enhanced geometric details, which is advantageous for applications like object detection and 3D reconstruction. The technique can efficiently capture surface details and fine-grained features by utilizing perturbation learning, which improves the overall quality of the unsampled point clouds. Furthermore, the method is adaptable and does not require specific surface representations, making it suitable for wide range of point cloud datasets. The computational complexity of training the perturbation model, which can rise in tandem with the volume and intricacy of the input point cloud, is one of the possible disadvantage, though. Additionally, the quality of the perturbation model that is trained and the availability of the adequate training data may determine how effective the method is, which may restrict its application to specific situations. Overall, \cite{ding2021point} presented encouraging advances in point cloud processing, nonetheless, for practical implementation, it is imperative to overcome the issues associated to computing and data. \cite{song2021layer} presented a novel method for geometric feature preservation in \ac{LiDAR} point cloud compression. This framework's ability to produce lossless compression-which guarantees that no information across many layers of the point cloud by utilizing a layer-wise geometry aggregation mechanism. This allows for efficient modelling and compression of intricate 3D structures. Furthermore, the framework is flexible and can be adjusted to suit various compression needs and \ac{LiDAR} datasets. A possible disadvantage is the computational complexity involved in analyzing and combining geometric data over several layers. This complexity may rise as input point cloud data becomes more complicated and large. Additionally, even though the technique produces lossless compression, it might not always outperform lossy compression techniques in terms of compression ratios, which might be important to take into account in situations when storage is limited. \cite{zhao2021transformer3d} offered a cutting-edge method for employing transformer networks to refine votes in order to improve 3D object detection performance. The method's main contribution is how it uses transformer architectures to combine and fine-tune votes from various point cloud regions in order to improve object detection. The approach improved accuracy and robustness in 3D object detection tasks by capturing global context information and spatial links between object proposals through the incorporation of transformer-based refinement. Furthermore, the approach is flexible and may be combined with current 3D object detection frameworks, offering a flexible way to improve detection efficiency. While these approaches were remarkably successful, in order to do feature learning, they need supervised data. Point cloud models are difficult to implement in novel real-world scenarios where labelled data is hard to come by due to their reliance on annotation. Therefore, its critical to devise strategies for lowering the quality of annotated samples needed to complete deep learning-based point cloud interpretation tasks with the necessary performance\cite{mei2022unsupervised}.